\subsection{Global integral completion with a reference localized at two}

\begin{theorem}[AT29: Proved; integral completion through form three with actual local reference paths]
\label{AT29}
Put $S=\Gamma_{\mathbb Z}(T)$.
Let $A_i$ be a specified finite nonempty family of torsionfree admissible
integral theories. Let $X$ be a strict simplicial PD cochain $S$-algebra
with torsionfree terms, coefficient $H^0$, and rational coefficient
comparison in every simplex and weight, with actual maps $f_i:X\to A_i$.
Let $D$ be a strict simplicial PD cochain $S_{(2)}$-algebra with
torsionfree terms, coefficient $H^0$, and rational coefficient comparison
in every simplex and weight. Specify actual maps
\[
 g:X_{(2)}\longrightarrow D,\qquad
 r_i:D\longrightarrow A_{i,(2)},\qquad
 F_i:X_{(2)}\longrightarrow A_{i,(2)}^{\Delta[1]},\qquad
 e_0F_i=r_i g,\quad e_1F_i=f_{i,(2)}.
 \tag{AT.200}
\]
Assume $D$ has coefficient comparison through form three for weights
$w\geq L$, uniformly in the simplex. Assume also that $X[1/2]$ has
that comparison for weights $w\geq L_o$, uniformly in the simplex.
Set
\[
 W=\max\{4,L,L_o,\lambda_{A_i}(6),\mu_{A_i}(6):i\}.
 \tag{AT.201}
\]
There are actual inclusions $X\hookrightarrow E$ over $\mathbb Z$
and $D\hookrightarrow D'$ over $\mathbb Z_{(2)}$, extending all the
maps in \textup{(AT.200)}, such that $E$ and $D'$ have torsionfree
terms, coefficient $H^0$, rational coefficient comparison in every
simplex and weight, and coefficient comparison through form three
for $w\geq W$, uniformly in the simplex. All the maps $E\to A_i$
are integral. All horizontal rows of both $E$ and $D'$ are eventually
acyclic, with the possible cutoff
\[
                  M(q)=\max\{q+1,\mu_{A_i}(q):i\}.
 \tag{AT.202}
\]
The support condition $q\leq w$ is retained when it holds in the
specified input. No comparison above form three or strict
identification of a low-form truncation is asserted.
\end{theorem}

\begin{proof}
All localizations and intersections below are taken inside the
specified rationalizations, into which the torsionfree terms embed.
Write
\[
 K_i=\ker\bigl(e_0:A_{i,(2)}^{\Delta[1]}\longrightarrow A_{i,(2)}\bigr).
 \tag{AT.203}
\]
For weights at least $W$, A-002 gives $H^j((K_i)_n)=0$
for $0\leq j\leq3$, in every simplex $n$.
The proof there uses only target row and vertical bounds through
degree six, included in \textup{(AT.201)}.
Its rows are contractible in every weight and form.
We also have the actual surjection
\[
                  e_1:(K_i)_n^{w,j}\longrightarrow
                          (A_{i,(2)})_n^{w,j}
       \quad (w\geq W,\ 0\leq j\leq2).
 \tag{AT.204}
\]
Indeed prescribe zero on $\Delta[n]\times\{0\}$ and the chosen
element on $\Delta[n]\times\{1\}$. A-001's finite restriction
epimorphism for the acyclic row in form $j$ extends those values
across $\Delta[n]\times\Delta[1]$.
This is a surjection of actual row modules, rather than an assertion
that a vertical primitive has a chosen endpoint.

\paragraph{Integral endpoints for a local two-cell extension.}
For $q=2$ or $3$, choose actual homogeneous data in the current source
\[
 z\in X_n^{w,q},\quad dz=0,\qquad
 b\in X_n^{w,q-1},\quad db=mz,
 \qquad w\geq W,\quad m=2^k,\quad k\geq1.
 \tag{AT.205}
\]
Choose $u_D,v_D\in D_n$, of forms $q-1,q-2$, with
$du_D=g(z)$ and $dv_D=mu_D-g(b)$.
They exist by the reference comparison in this weight.
The endpoint-zero-kernel argument of AT27 gives local path values
$U_{i,\mathrm{loc}},V_{i,\mathrm{loc}}$ with
\[
 \begin{aligned}
 dU_{i,\mathrm{loc}}&=F_i(z),&
 dV_{i,\mathrm{loc}}&=mU_{i,\mathrm{loc}}-F_i(b),\\
 e_0U_{i,\mathrm{loc}}&=r_i(u_D),&
 e_0V_{i,\mathrm{loc}}&=r_i(v_D).
 \end{aligned}
 \tag{AT.206}
\]
Explicitly subtract the constant reference value, fill its closed
error in $K_i$, and repeat for the second equation. This uses only
cohomology in forms $q,q-1$ of $K_i$.

For every abelian group $M$, the actual map
$M/mM\to M_{(2)}/mM_{(2)}$ is an isomorphism: all odd integers
act invertibly on modules killed by $m=2^k$.
Apply this degreewise to the cochain complex $(A_i)_n$.
The endpoint-one pair in \textup{(AT.206)} shows that $f_i(b)$
modulo $m$ is a boundary. Choose $v_i\in(A_i)_n^{w,q-2}$
with $dv_i\equiv f_i(b)\pmod m$ and put
\[
             U_i=\frac{f_i(b)-dv_i}{m},\qquad V_i=-v_i.
 \tag{AT.207}
\]
These are actual integral elements. Torsionfreeness and $db=mz$
give $dU_i=f_i(z)$; directly $dV_i=mU_i-f_i(b)$.

We now impose these integral endpoints on the paths.
Put $\delta U=e_1U_{i,\mathrm{loc}}-U_i$ and
$\delta V=e_1V_{i,\mathrm{loc}}-V_i$.
They satisfy $d\delta U=0$ and $d\delta V=m\delta U$.
Target vertical exactness gives $K\in(A_{i,(2)})_n^{w,q-2}$
with $dK=\delta U$.
For $q=3$, the closed element $\delta V-mK$ has form one;
choose $H$ of form zero with $dH=\delta V-mK$.
For $q=2$, that closed element is a coefficient $c\in S_{(2)}^w$.
Choose $a\in S^w$ whose image modulo $mS_{(2)}^w$ is $c$.
Replace $V_i$ by $V_i+a$, and replace $K$ by $K+(c-a)/m$.
The latter is a coefficient correction, so $dK$ is unchanged,
and now $\delta V-mK=0$ exactly. In this case put $H=0$ in
degree minus one. These formulas fix the sign and the integral
representative in the degree-zero endpoint correction.

Lift $K,H$ through \textup{(AT.204)} to
$\widetilde K,\widetilde H$ in $K_i$, using zero for the
degree-minus-one term. Replace the paths by
\[
 U'_{i,\mathrm{loc}}=U_{i,\mathrm{loc}}-d\widetilde K,
 \qquad
 V'_{i,\mathrm{loc}}=V_{i,\mathrm{loc}}
                         -m\widetilde K-d\widetilde H.
 \tag{AT.208}
\]
Their differential equations and endpoint-zero values are unchanged;
their endpoint-one values are exactly the integral $U_i,V_i$.

Adjoin free simplicial $\mathbb Z[\Delta[n]]$ copies of generators
$u,v$ of forms $q-1,q-2$ and weight $w$, with $du=z$, $dv=mu-b$.
Representability makes these data and the chosen values actual
simplicial parameter maps, including every face and degeneracy.
Send $u,v$ to $U_i,V_i$ in the integral targets, to $u_D,v_D$
in the local reference, and to \textup{(AT.208)} in the local
path objects. The relative free PD universal property gives the
strict maps and both endpoint identities.
Rationally the actual change $u'=u-b/m$ gives $du'=0$, $dv=mu'$,
a contractible free graded-commutative disk.
Take the image of the extension in its rationalization.
The old source remains injective; every receiving object is
torsionfree, so all of the specified maps factor through this image.

\paragraph{The comparison after inverting two.}
Over $\mathbb Z[1/2]$ the same change $u'=u-b/m$ is integral.
For $q=2$ it gives a free even-zero-to-odd-one identity disk,
which contracts integrally. For $q=3$ it gives a free
odd-one-to-even-two identity disk. Its arity-$a$ pair has
differential multiplication by $a$, in forms $2a-1,2a$.
The arities $a=1,2$ contract because $1,2$ are units, and
every remaining pair begins in form at least five.
Thus tensoring with these disks preserves cohomology through
form three by actual linear contractions; no free-PD comparison
on a nonfree module is used. The representable parameter is free
in each simplex. Finite products and filtered unions therefore
give the result for all new generators.
The displayed localized free-disk algebra is torsionfree, so it
is also the localization of the chosen torsionfree image.

Consequently the comparison for $X[1/2]$ through form three and
its uniform bound $L_o$ persist at every such stage.
Exactness of localization shows that every individual positive
cohomology class in this range is killed by some power of two.
This supplies data \textup{(AT.205)} for every class in forms
two and three, without a finiteness assumption on a whole
cohomology module.

\paragraph{Saturation with dyadic denominators and a local reference.}
For a current integral source $Z$, inside $Z[1/2]$ adjoin all
homogeneous zero-forms $a$ of weights at least $W$ such that
\[
                         da\in Z,
                  \qquad g_{\mathbb Q}(a)\in D.
 \tag{AT.209}
\]
The conditions are stable under sums, coefficient multiples and
simplicial structure maps. Take their PD algebra and repeat
countably. The differential is defined because each new derivative
was integral; all existing relations hold in the fixed rational
algebra. Rationalization does not change.

The local paths extend integrally over $\mathbb Z_{(2)}$.
For a root $a$, subtract the constant path on $r_i g_{\mathbb Q}(a)$
from $F_{i,\mathbb Q}(a)$. The difference lies in $(K_i)_{\mathbb Q}^0$
and its differential is integral in $(K_i)^1$.
The vanishing of $H^1(K_i)$ gives an integral primitive there.
The difference from it is a rational closed zero-form of $K_i$,
which is zero by $H^0((K_i)_{\mathbb Q})=0$.
Thus the path value, and in particular its endpoint-one value,
belong to the local integral objects.

That endpoint value $f_{i,\mathbb Q}(a)$ also belongs to $A_i[1/2]$.
For every torsionfree abelian group $M$ one has, inside $M_{\mathbb Q}$,
\[
                         M[1/2]\cap M_{(2)}=M.
 \tag{AT.210}
\]
Indeed write $x=b/2^k=c/s$, where $s$ is odd, so $sb=2^kc$.
Choose integers $u,v$ with $u2^k+vs=1$; then $x=ub+vc\in M$.
Therefore every endpoint image is actually integral in $A_i$.
This extends the original integral maps, and their localizations
retain the strict endpoint identity with the paths.

This saturation kills all closed one-forms in the stated range.
For such a form $t$, the comparison after inverting two gives
$db=2^kt$ for an integral zero-form $b$. Put $a_0=b/2^k$.
Choose a local reference primitive $b_D$ of $g(t)$.
Then $g_{\mathbb Q}(a_0)-b_D$ is a rational coefficient $c$.
The coefficient identity
$\mathbb Q=\mathbb Z[1/2]+\mathbb Z_{(2)}$ lets us choose
$c_0\in S_{\mathbb Z[1/2]}^w$ with $c-c_0\in S_{(2)}^w$.
Explicitly, for $a/(2^ks)$ with $s$ odd choose $b$ with
$bs\equiv a\pmod {2^k}$ and subtract $b/2^k$; the remainder
has odd denominator. Thus $a=a_0-c_0$ satisfies
\textup{(AT.209)} and $da=t$.

Saturation leaves the algebra after inverting two literally
unchanged, because every adjoined root already lies in $Z[1/2]$.
It also preserves coefficient $H^0$: a new closed zero-form is a
rational coefficient in both $S_{\mathbb Z[1/2]}$ and $S_{(2)}$,
hence belongs to $S$ by \textup{(AT.210)}.
The first membership uses $H^0(Z[1/2])=S_{\mathbb Z[1/2]}$;
the second uses its image in the reference.
After a two-cell extension the same coefficient assertion follows
in its new weights from the actual integral map to any target,
where $H^0=S$ in that range; below the new weights it is unchanged.
Torsionfreeness and rational comparison hold throughout.

Alternate the saturation with simultaneous extensions for all
nonzero classes in forms two and three, every simplex, and every
weight at least $W$. Each stage has a set of classes and actual
representatives. Every element in the compatible countable union
belongs to a stage; each positive closed form is killed at a later
stage. This gives global comparison through form three with the
uniform cutoff $W$, retaining all integral maps and local paths.

\paragraph{Horizontal restoration with both endpoints fixed.}
Restore the rows simultaneously in the integral source and local
reference. For a source cycle $z\in N_nE^{w,q}$, $w\geq M(q)$,
adjoin a reference C21 free disk with normalized top $a_D$ in
simplex $n+1$, form $q$, and boundary $g(z)$.
For each target choose a local normalized filler $b_i$ of
$r_i g(z)$, and an integral normalized filler $a_i$ of $f_i(z)$.
They exist by row acyclicity in \textup{(AT.202)}.
Send $a_D$ to $b_i$, and adjoin the corresponding integral source
disk top $a$ with boundary $z$, mapped to $a_D$ in the reference
and $a_i$ in the integral target.

The new path is an explicit finite restriction problem in the
local target row. On
\[
 (\partial\Delta[n+1]\times\Delta[1])
 \ \cup\ (\Delta[n+1]\times\partial\Delta[1])
 \ \subset\ \Delta[n+1]\times\Delta[1]
\]
give face zero the prescribed path $F_i(z)$, every other lateral
face zero, endpoint zero the simplex $b_i$, and endpoint one the
simplex $a_i$. All intersections agree: the zero face at the
endpoints is respectively $r_i g(z)$ and $f_i(z)$, and the other
faces vanish by normalization.
A-001 extends this boundary across the product at the stated
weight. The resulting normalized path is the new top value;
the differential partner is sent to its differential.
Thus both endpoints and every face equation hold strictly.

Also adjoin reference-only disks for its other cycles in this
range, choosing local target fillers. Existing source maps compose
with the reference inclusion and the paths stay fixed.
Iterating these two operations kills every row cycle at the
asserted form-dependent cutoff.
In every simplex, a relative C21 disk has free identity-disk
complements. The direct arity-pair argument in AT27 applies over
$\mathbb Z$ and over $\mathbb Z_{(2)}$: the first nonunit
odd-disk pair starts in form three, where its tensor's cycles
are zero because multiplication by its integer is injective on
coefficient $H^0$. Later pairs start above the range; even disks
contract. Thus these extensions preserve cohomology through form
three even for torsionfree current terms. Finite products and
filtered unions preserve this argument weightwise.
They preserve $H^0$, rational comparison and torsionfreeness too.
The bound $M(q)\geq q+1$ handles the support of each differential
partner; \textup{(AT.201)} handles all earlier two-cell and root
additions. This proves every assertion.
\end{proof}

\begin{corollary}[AT29b: Proved; an actual global free-family application]
\label{AT29b}
Take the integral free-cell data and weight bounds of A-S09-31,
and impose the actual global relative free $D(2)$ presentation
and free identity-disk base $P$ of AT26.
For the initial integral resolved extension $Y$ there are actual
inclusions $Y\hookrightarrow E$ and integral maps $E\to A,B$
extending the specified maps on the old source, with the global
coefficient comparison through form three and all eventual
horizontal row comparisons of AT29.
Every simplex degree is included in that vertical assertion.
\end{corollary}

\begin{proof}
Use the initial integral lift choices in A-S09-31, before its
subsequent simplex-two cleanup. Its formulas (S09.170)--(S09.172)
align each localized carrier with the chosen integral map on the
same integral coefficient resolution. The correction has every
positive face zero and leaves the reference values unchanged.
Consequently the complete prism in AT26 has the same two other
tetrahedra and the same constant bottom prism; its corrected first
tetrahedron has the required integral endpoint. The complete
face-table proof of AT26 gives the maps \textup{(AT.200)} with
its actual common reference $D$.
That reference has coefficient comparison through form three
in every simplex and weight, as proved there.

We check the source after inverting two in every simplex degree.
The global free presentation writes the old source simplexwise
as $S$ tensored with free identity disks, including the complements
of its relative selected disks. Over $\mathbb Z[1/2]$ the exact
graph and inverse changes (S09.173)--(S09.175) apply to every
resolved summand and every basis value of its free $K(2)$ parameter.
They turn the coefficient module, including its formerly nonfree
bottom, into free disks. The affine twist is removed by the actual
change $u_0'=u_0-\varepsilon\eta/2$.
Even disks contract. In each odd-one-form disk its arities one
and two contract, and every remaining arity starts in form at
least five. Higher odd-form disks start at least as late.
Therefore $Y[1/2]$ has coefficient comparison through form three
in every simplex and weight. This localized free presentation is
torsionfree, so it is exactly the localization of the integral
torsionfree image, by the same argument as in A-S09-31.
All other input and map properties were proved for this initial
extension in AT16 and A-S09-31. Applying AT29 proves the claim.
\end{proof}

These inclusions can add forms zero and one. They do not supply
the strict low-form isomorphisms in C34. AT28 gives an actual
coherent-reference instance showing why those isomorphisms cannot
be inferred from reference-path data alone. No comparison in form
four or above is included in these conclusions.
